\begingroup
\renewcommand{\arraystretch}{1.45}
\Needspace{46\baselineskip}
\InterpolationArticle{29}
Es bleiben die Fälle, in denen in der Funktion $X$ nur Kosinus oder nur Sinus vorkommen, die wir in den Artikeln 16, 17 allgemein behandelt haben. Wie zuvor werden wir voraussetzen, dass eine Periode von $\mu$ Gliedern gegeben ist, etwa $X=A,B,C,D$ usw. für $x=a,b,c,d$ usw. Hier sind jedoch zwei Fälle sorgfältig zu unterscheiden, da sie verschieden behandelt werden müssen:

der $1^{\mathrm{te}}$ Fall, in dem $a$ das Komplement eines der folgenden Werte $b,c,d$ usw. zu $360^{0}$ oder zu $0$ oder zu einem Vielfachen des Umfangs ist. Dieses Glied sei $a+\frac{k}{\mu}360^{0}$,
\SourceLeaf{320}
und es sei $a+\frac{k}{\mu}360^{0}+a=l\times360^{0}$ oder $a=(\mu l-k)\frac{180}{\mu}$. Daher wird in diesem Fall $a$ durch $\frac{180}{\mu}$ teilbar sein, also $\sin\mu a=0$; ferner gilt $b=(\mu l-k+2)\frac{180^{0}}{\mu}$ und $\sin\mu b=0$, und ebenso für die folgenden Werte $c,d$ usw. Unter der Voraussetzung also, dass $a$ kleiner als $\frac{360^{0}}{\mu}$, aber nicht negativ ist (was zulässig ist, da man von allen $a,b,c,d$ usw. das nächstkleinere Vielfache des ganzen Umfangs abziehen und den kleinsten der Reste, der jene Eigenschaft haben wird, als Anfang der Periode nehmen darf), wird entweder $a=0$ und $b=\frac1\mu360^{0}$, $c=\frac2\mu360^{0}$, $d=\frac3\mu360^{0}$ usw. gelten; oder $a=\frac1\mu180^{0}$, $b=\frac3\mu180^{0}$, $c=\frac5\mu180^{0}$, $d=\frac7\mu180^{0}$ usw.

\Needspace{22\baselineskip}
und der $2^{\mathrm{te}}$ Fall, in dem es nicht das Komplement irgendeines der folgenden Werte $b,c,d$ usw. zu $360^{0}$ oder $0$ oder einem Vielfachen des Umfangs ist. Hier wird $a$ nicht durch $\frac{180^{0}}{\mu}$ teilbar sein, also $\sin\mu a=\sin\mu b=\sin\mu c=\sin\mu d$ usw. nicht $=0$.

Im letzteren Fall kann die allgemeine Methode des Artikels 16 sofort angewandt werden. Im ersteren müssen dagegen von den Werten $a,b,c,d$ usw. (die wir als unter $360^{0}$ reduziert voraussetzen) zuvor alle verworfen werden, die größer als $180^{0}$ sind, da sich ihre Komplemente zu $360^{0}$ unter den übrigen finden; außerdem, wenn sie vorkommen, die Werte $0$ und $180^{0}$, sofern in der Funktion $X$ nur Sinus auftreten. Da die Anwendung der allgemeinen Methode auf diese Weise jedoch weniger zweckmäßig wäre, werden wir diesen Fall unten auf andere Weise erledigen. Wir werden nun mit dem Fall beginnen, in dem $\sin\mu a$ nicht $=0$ ist.
\Needspace{37\baselineskip}
\InterpolationArticle{30}
Zunächst sei $X$ eine Funktion der Form
\[
\begin{aligned}
&\alpha+\alpha'\cos x+\alpha''\cos2x+\alpha'''\cos3x+.\,.\,.+\alpha^{n}\cos nx
\end{aligned}
\]
somit muss $\mu=n+1$ sein. Hier gilt nach dem zweiten Lemma des Artikels 19
\[
\begin{aligned}
&(\cos t-\cos b)(\cos t-\cos c)(\cos t-\cos d)\ \text{usw.}
\\[5pt]
&=\frac1{2^{\mu-1}}\times\frac{\cos\mu t-\cos\mu a}{\cos t-\cos a}\\[5pt]
&=\frac1{2^{\mu-1}\sin a}\{\sin\mu a+2\sin(\mu-1)a\cos t\\[5pt]
&\qquad{}+2\sin(\mu-2)a\cos2t+\text{usw.}+2\sin a\cos(\mu-1)t\}
\end{aligned}
\]
Es lässt sich nämlich leicht bestätigen, dass das Produkt aus der Summe im zweiten Teil dieser Gleichung und $\cos t-\cos a$ den Wert $=\sin a(\cos\mu t-\cos\mu a)$ erhält. Daher ergibt sich auch, wenn man $a$ für $t$ schreibt,
\[
\begin{aligned}
&(\cos a-\cos b)(\cos a-\cos c)(\cos a-\cos d)\ \text{usw.}
\\[5pt]
&=\frac1{2^{\mu-1}\sin a}\{\sin\mu a+2\sin(\mu-1)a\cos a\\[5pt]
&\qquad+2\sin(\mu-2)a\cos2a+\text{usw.}+2\sin a\cos(\mu-1)a\}
\end{aligned}
\]
\Needspace{14\baselineskip}
Die Summe in diesem Ausdruck wird $=\mu\sin\mu a$, wie daraus ersichtlich ist, dass
\SourceLeaf{321}
\[
\begin{aligned}
&\text{zweites Glied}=\sin\mu a+\sin(\mu-2)a\\[5pt]
&\text{letztes Glied}=\sin\mu a-\sin(\mu-2)a\\[5pt]
&\text{ferner, drittes Glied}=\sin\mu a+\sin(\mu-4)a\\[5pt]
&\text{vorletztes Glied}=\sin\mu a-\sin(\mu-4)a\\[5pt]
&\text{usw.}
\end{aligned}
\]
\Needspace{46\baselineskip}
Hieraus folgt, dass der Koeffizient von $A$ in der Formel des Artikels 16 für $T$ gleich
\[
\begin{aligned}
&\frac1{\mu\sin\mu a}\{\sin\mu a+2\sin(\mu-1)a\cos t\\[5pt]
&\qquad{}+2\sin(\mu-2)a\cos2t+\text{usw.}+2\sin a\sin(\mu-1)t\}
\end{aligned}
\]
wird. Ganz ähnliche Ausdrücke ergeben sich für die Koeffizienten von $B,C,D$ usw., indem man lediglich $a$ in $b,c,d$ usw. ändert. Setzt man daher
\[
\begin{aligned}
&\varepsilon=\frac1\mu(A+B+C+D+\text{usw.})\\[5pt]
&\varepsilon'=\frac2{\mu\sin\mu a}\{A\sin(\mu-1)a+B\sin(\mu-1)b\\[5pt]
&\qquad{}+C\sin(\mu-1)c+D\sin(\mu-1)d+\text{usw.}\}\\[5pt]
&\varepsilon''=\frac2{\mu\sin\mu a}\{A\sin(\mu-2)a+B\sin(\mu-2)b\\[5pt]
&\qquad{}+C\sin(\mu-2)c+D\sin(\mu-2)d+\text{usw.}\}\\[5pt]
&\varepsilon'''=\frac2{\mu\sin\mu a}\{A\sin(\mu-3)a+B\sin(\mu-3)b\\[5pt]
&\qquad{}+C\sin(\mu-3)c+D\sin(\mu-3)d+\text{usw.}\}\\[5pt]
&\text{usw.}\\[5pt]
&\varepsilon^{n}=\frac2{\mu\sin\mu a}(A\sin a+B\sin b\\[5pt]
&\qquad{}+C\sin c+D\sin d+\text{usw.})
\end{aligned}
\]
so wird
\[
\begin{aligned}
&T=\varepsilon+\varepsilon'\cos t+\varepsilon''\cos2t+\varepsilon'''\cos3t+\text{usw.}\\[5pt]
&\qquad{}+\varepsilon^{n}\cos nt
\end{aligned}
\]
Da diese Formel allgemein für jeden beliebigen Wert von $t$ gelten muss, wird sie notwendig mit $X$ identisch werden, wenn man $t$ in $x$ ändert, und somit werden die Koeffizienten $\varepsilon,\varepsilon',\varepsilon''$ usw. jeweils $\alpha',\alpha',\alpha''$ usw. gleich sein.
\endgroup
